\section*{ID7794: Coupled Action / Energy Functional: α(T)}
\paragraph{Metadata.}
PiID: 8400269104140 | RTEID: RTE:P30265.8329:T5.8865 | UUID: 54619037-b333-5edb-bc81-e4fafed79792

\paragraph{Canonical Statement.}
Take free energy (Ginzburg–Landau) with acoustic drive: \textbackslash{}[ \textbackslash{}mathcal\{F\}[\textbackslash{}Psi,A,p]=\textbackslash{}int d\textasciicircum{}3x\textbackslash{};\textbackslash{}Bigg\textbackslash{}\{ \textbackslash{}alpha(T)|\textbackslash{}Psi|\textasciicircum{}2 + \textbackslash{}frac\{\textbackslash{}beta\}\{2\}|\textbackslash{}Psi|\textasciicircum{}4 + \textbackslash{}frac\{1\}\{2m\textasciicircum{}*\}\textbackslash{}left|\textbackslash{}left(-i\textbackslash{}hbar\textbackslash{}nabla - q\textasciicircum{}* \textbackslash{}mathbf\{A\}\textbackslash{}right)\textbackslash{}Psi\textbackslash{}right|\textasciicircum{}2 + \textbackslash{}frac\{|\textbackslash{}mathbf\{B\}|\textasciicircum{}2\}\{2\textbackslash{}mu\_0\} + \textbackslash{}frac\{p\textasciicircum{}2\}\{2\textbackslash{}rho\} + \textbackslash{}frac\{\textbackslash{}rho c\_s\textasciicircum{}2\}\{2\}|\textbackslash{}nabla \textbackslash{}Phi\_a|\textasciicircum{}2 - g\textbackslash{}, p(\textbackslash{}mathbf\{x\},t)\textbackslash{},|\textbackslash{}Psi|\textasciicircum{}2 \textbackslash{}Bigg\textbackslash{}\}, \textbackslash{}] where - \textbackslash{}(m\textasciicircum{}*\textbackslash{}), \textbackslash{}(q\textasciicircum{}*=2e\textbackslash{}) are Cooper pair mass and charge; - \textbackslash{}(\textbackslash{}alpha(T)=\textbackslash{}alpha\_0(T-T\_c)\textbackslash{}), \textbackslash{}(\textbackslash{}beta>0\textbackslash{}); - \textbackslash{}(p\textbackslash{}) is acoustic pressure, \textbackslash{}(\textbackslash{}Phi\_a\textbackslash{}) acoustic potential with \textbackslash{}(p=\textbackslash{}rho\textbackslash{}partial\_t \textbackslash{}Phi\_a\textbackslash{}) or simply treat driven pressure \textbackslash{}(p\textbackslash{}); - \textbackslash{}(g\textbackslash{}) is the coupling coefficient (acoustic modulation of pair condensation energy).

\paragraph{Canonical Equation.}
\[
\alpha(T)=\alpha_0(T-T_c)
\]

\paragraph{Dictionary.}
Coupled Action / Energy Functional: α(T) — α(T)=α\_0(T-T\_c)

\paragraph{Encyclopedia.}
Coupled Action / Energy Functional: α(T) is an algebraic definition, written α(T)=α\_0(T-T\_c). It is an algebraic definition expressing the quantity directly in terms of the variables on its right-hand side. It is built from T, T\_c, defining α(T). Within the Trawin Zero Point registry it serves as one element of the framework's mathematical narrative.
